\cotefile{157-3}
\scanpage{252}
\mspage{202}

\[
(\unsure{\xi|\zeta}*\xi')*\xi''\ \unsure{\xrightarrow{\sim}|\overset{\sim}{\text{---}}}\ \xi*(\xi'*\xi'')
\]
Il \deleted{\illegible{1}} y a deux façons d'\unsure{immerger|envoyer} un produit binaire
\[
\xi*\xi',\quad \xi'*\xi'',\quad \xi*\xi''
\]
dans le produit ternaire, \unsure{eh|et} bien elles doivent \deleted{\unsure{dans}} chacun des trois cas, être égales. Et de \unsure{$=$|même} il y a, ceci admis, \deleted{\unsure{deux une}} deux façons d'\unsure{immerger|envoyer} $\xi$ dans le produit ternaire, en factorisant soit par $\xi*\xi'$, soit par $\xi*\xi''$, et de $=$ pour les deux autres facteurs $\xi',\xi''$; et là encore, ces deux façons doivent être égales, dans chacun des \unsure{ces|en} trois cas -- donc $3+3=6$ égalités en \unsure{tout} à écrire \inserted{et à vérifier} -- on a vu pire !

La philosophie \unsure{calculatoire} devrait être celle-ci~: \unsure{Si on a|Lié à} un produit $n$-uple
\[
\xi_1*\xi_2*\text{---}*\xi_n
\]
\unsure{(ici|(où} je me \unsure{dispense} de mettre des \unsure{parenthèses} ....), pour \unsure{tout} produit $m$-uple partiel ($m\in[1,n-1]$)
% archive sheet number: 252
